\begin{titlepage}
  \centering
  \vspace*{1.5cm}
  {\Large\textsf{\manualeditor}\par}
  \vspace{1.2cm}
  {\fontsize{36}{42}\selectfont\bfseries\color{vrpdark}\manueltitle\par}
  \vspace{0.6cm}
  {\LARGE\manualsubtitle\par}
  \vspace{1.5cm}
  \begin{tikzpicture}
    \fill[vrpblue, rounded corners=5pt] (0,0) rectangle (11,1);
    \node[white, font=\small\bfseries] at (5.5,0.5)
      {Structures · Prestations · Séances · Patients · Facturation · Trésorerie};
  \end{tikzpicture}
  \vfill
  {\large\manualversion\par}
  \vspace{0.4cm}
  {\small Guide court — profil medical\par}
\end{titlepage}

\section*{À propos de ce guide}
\addcontentsline{toc}{section}{À propos de ce guide}

Ce document est un \textbf{guide de prise en main rapide} pour les professionnels de santé et structures de soin utilisant VRP Plan avec le \textbf{contexte métier médical}. Il ne couvre pas les profils formation ni conseil qui sont couverts dans d'autres guides.

\textbf{Durée estimée}~: 15 à 20 minutes pour la première configuration, puis quelques minutes par jour pour le suivi.

%\begin{astuce}[Vocabulaire]
%  Dans l'application, vous verrez \emph{Structure}, \emph{Prestation}, \emph{Séance} et \emph{Patient} à la place d'école, programme, cours et groupe. Le fonctionnement reste identique.
%\end{astuce}
